\section{Positive time: a growing clock window}\label{sec:positive}
The positive-time argument has three steps: control the complete source,
compare its transform with a carrier, and transfer zeros to a regularized
averaged density. All source indices are retained.

\subsection{A generic regularization bridge}
For a nonzero holomorphic function $H$ on a neighborhood of a compact
set, let $\nu=\sum_a m_a\delta_a$ be its zero measure.
The distributional identity $\Delta\log|H|^2=4\pi\nu$ gives
\begin{equation}
 \rh-\nu=\frac1{4\pi}\Delta
       \log\left(1+\frac{\epsilon^2}{|H|^2}\right).
 \label{eq:bridge}
\end{equation}
For a compactly supported $C^2$ test kernel $K$,
\begin{equation}
 |\langle\rh-\nu,K\rangle|
 \le\frac{\|\Delta K\|_\infty}{4\pi}
      \int_{\supp K}\log(1+\epsilon^2/|H|^2)\,dA .
 \label{eq:bridgebound}
\end{equation}
If $|H|\ge b$ outside zero disks, and inside a disk of radius $r_a$
one has $|H(z)|\ge c_a|z-a|^{m_a}$, its logarithmic integral is bounded
by
\[
 I_m(a,r)=2\pi\int_0^r s\log(1+(a/s)^{2m})\,ds,\qquad
 a=(\epsilon/c_a)^{1/m_a}.
\]
For simple zeros,
\begin{align}
 I_1(a,r)&=\pi\{r^2\log(1+a^2/r^2)
                 +a^2\log(1+r^2/a^2)\}\notag\\
 &\le \pi a^2[1+\log(1+r^2/a^2)] =:\overline I_1(a,r).
 \label{eq:I1}
\end{align}
For $m\ge2$, splitting the integral at $s=a$ yields
$I_m(a,r)\le\pi a^2[m+\log2+(m-1)^{-1}]$.
Thus simple and multiple zeros have distinct regularization rates.
The bridge does not exclude multiple zeros.

\subsection{Complete-source concentration at the large saddle}
Factor the source as
\begin{align}
 e^{tu^2}\Phi(u)&=2\pi^2e^{q_t(u)}b(u),&
 q_t(u)&=tu^2+9u-\pi e^{4u},\label{eq:q}\\
 b(u)&=\sum_{n\ge1}
  \left(n^4-\frac{3n^2e^{-4u}}{2\pi}\right)
  e^{-\pi(n^2-1)e^{4u}},&
 0<b(u)&\le C_b=(1-16e^{-3\pi})^{-1}.\notag
\end{align}
Recall the large saddle $L$ and curvature $\kappa$ from
\eqref{eq:saddle}.
It has the explicit representation and expansion
\begin{align}
 L&=-\tfrac14W_{-1}\left(-\frac{8\pi}{t}e^{-18/t}\right)
       -\frac9{2t},\label{eq:Lambert}\\
 L&=\tfrac14\{\log t+\log\log t-\log(8\pi)\}
       +O(\log\log t/\log t),\qquad \kappa\sim8tL .\notag
\end{align}
Define the probability law
\begin{gather*}
 M=H(t,0),\qquad dP_t(u)=M^{-1}e^{tu^2}\Phi(u)\,du,\\
 \mu=\E_tu,\qquad v=\Var_tu,\qquad
 B_3(Y)=\E_t[|u-\mu|^3e^{Y|u-\mu|}].
\end{gather*}
The finite quantities in Appendix~\ref{app:source} satisfy
\begin{equation}
 |\mu-L|\le e_\mu,\qquad v\le A_2(0),\qquad
 B_3(Y)\le\widehat B_3(Y),\qquad
 M\ge \frac{2\pi^2e^{q_t(L)}}{\sqrt\kappa}Z_- .
 \label{eq:momentbounds}
\end{equation}
For fixed $Y$, their orders are
\begin{equation}
 e_\mu=O(\kappa^{-3/4}),\quad
 v=O(\kappa^{-1}),\quad
 B_3(Y)=O(\kappa^{-3/2}),\qquad Z_-\longrightarrow\sqrt{2\pi}.
 \label{eq:momentorders}
\end{equation}
The signed mean error, not merely an absolute first moment,
will be essential for a growing Vandermonde.

\subsection{A finite carrier comparison}
Choose an integer $N\ge1$, put $X=N\pi/\mu$, and fix $Y>0$.
Let $Z=(X^2+Y^2)^{1/2}$ and define
\begin{equation}
 \begin{gathered}
 A(z)=1-\tfrac12vz^2,\qquad
 a_*=1-\tfrac12vZ^2>0,\qquad
 \delta=\frac{Z^3B_3(Y)}{6a_*}>0,\\
 s=4\delta\le\tfrac12,\qquad
 r=\frac{s}{\mu}<\min\{Y,\pi/(2\mu)\}.
 \end{gathered}
 \label{eq:carriergates}
\end{equation}
One can replace $v$ and $B_3(Y)$ in the bounds by their explicit
upper allowances, keeping the actual $v$ in $A(z)$.

\begin{lemma}[Carrier zeros and lower bounds]\label{lem:carrier}
Under \eqref{eq:carriergates}, on $|\Re z|\le X$, $|\Im z|\le Y$,
\begin{equation}
 \left|\frac{H(t,z)}{MA(z)}-\cos(\mu z)\right|
       \le\delta e^{\mu|\Im z|}.
 \label{eq:carrier}
\end{equation}
There is exactly one simple real zero $a_j$ in each disk of radius $r$
about $c_j=(j+\tfrac12)\pi/\mu$ in this rectangle, and no other zeros.
Furthermore
\begin{align}
 |a_j-c_j|&\le r,\qquad
 |H/M|\ge b_*:=a_*(e^{-s}\sin s-\delta)\ge a_*\delta
              &&\text{outside the disks},\label{eq:floor}\\
 |H(z)/M|&\ge c_*|z-a_j|,\qquad c_*=\frac{b_*}{2r}
              &&\text{inside the $j$th disk}.\label{eq:quotient}
\end{align}
\end{lemma}
\begin{proof}
The centered characteristic transforms obey
$\E e^{\pm iz(u-\mu)}=1-vz^2/2+\mathcal R_\pm$ with
$|\mathcal R_\pm|\le |z|^3B_3(Y)/6$.
Writing the cosine as two exponentials proves \eqref{eq:carrier}.
Outside the disks,
\begin{equation}
 |\cos(\mu z)|\ge e^{\mu|\Im z|-s}\sin s .
 \label{eq:cosfloor}
\end{equation}
For completeness, let $a$ be the distance of $\mu\Re z$ from the
nearest real cosine zero, reduced to $[0,\pi/2]$, and
$b=\mu|\Im z|$. If $b\ge s$, use
$e^{-b}\sinh b\ge(1-e^{-2s})/2\ge e^{-s}\sin s$.
If $b\le s$ and $a^2+b^2\ge s^2$, then
$\sin^2a+\sinh^2b\ge\sin^2s$: for $a\le s$ use the decrease
of $\sin a/a$, and for $a\ge s$ use monotonicity of sine.
This proves \eqref{eq:cosfloor}.
Since $e^{-s}\sin s\ge s/2=2\delta$ for $s\le1/2$,
Rouche's theorem \cite{DLMF110Coarse2026} gives one zero in each disk.
Conjugation makes it real, and multiplicity one follows from the count.
The same comparison excludes zeros elsewhere.
The quotient $g(z)=(H/M)/(z-a_j)$ is nonvanishing in the disk;
its boundary modulus is at least $b_*/(2r)$.
Applying the maximum-modulus principle to $1/g$ proves
\eqref{eq:quotient}.
\end{proof}

\subsection{A local coarse-density error}
Use the fixed spatial kernel \eqref{eq:kernel}.
Suppose the observation center satisfies
\begin{equation}
 |x|+\ell+r<X,\qquad |y|+h\le Y .
 \label{eq:buffer}
\end{equation}
Only the carrier disks with $|c_j-x|\le\ell+r$ can affect this
observation; their number is at most
$n_\ell=2\mu(\ell+r)/\pi+2$.
Put $\eta=\epsilon/M$ and
\begin{align}
 E_{\rm ph}&=\frac{(\pi/\mu)^3}{720}\|k_\ell^{(4)}\|_1,\label{eq:Ephase}\\
 E_{\rm reg}&=\frac{\|\Delta K\|_\infty}{4\pi}
 \left\{\frac{4\ell h\,\eta^2}{b_*^2}
              +n_\ell\overline I_1(\eta/c_*,2r)\right\}.\label{eq:Ereg}
\end{align}

\begin{theorem}[Finite physical-density allowance]\label{thm:positivefinite}
Under the source and carrier gates,
\begin{equation}
 \left|K*\rh(t,x,y)-\frac{L(t)}{\pi}k_h(y)\right|
 \le k_h(y)\left\{\frac{e_\mu}{\pi}
       +n_\ell r\|k_\ell'\|_\infty+E_{\rm ph}\right\}+E_{\rm reg}.
 \label{eq:positivefinite}
\end{equation}
\end{theorem}
\begin{proof}
Match each actual zero to its carrier disk. The resulting kernel
difference is at most $r\|k_\ell'\|_\infty k_h(y)$ per contributing
disk; use the local count $n_\ell$, not the full zero population.
For lattice spacing $a=\pi/\mu$, periodization and four integrations
by parts give
\[
 \left|\sum_jk_\ell(x-(j+\tfrac12)a)-a^{-1}\right|
 \le \frac{2\zeta(4)a^3}{(2\pi)^4}\|k_\ell^{(4)}\|_1=E_{\rm ph}.
\]
Replacing $\mu$ by $L$ costs $e_\mu k_h(y)/\pi$.
Finally apply \eqref{eq:bridgebound} to $H/M$ with regulator $\eta$.
The outside part has area at most $4\ell h$ and floor $b_*$.
Each intersecting carrier disk lies within distance $2r$ of its
actual zero, where \eqref{eq:quotient} and \eqref{eq:I1} apply.
This gives \eqref{eq:Ereg}.
\end{proof}

With $C=315/256$, explicit usable kernel constants are
\[
 \|k_\ell'\|_\infty=\frac{8C(6/7)^3}{\sqrt7\,\ell^2},\qquad
 \|\Delta K\|_\infty\le8C^2(\ell^{-3}h^{-1}+\ell^{-1}h^{-3}).
\]
Also $\|k_\ell^{(4)}\|_1=\ell^{-4}\|k^{(4)}\|_1$, and inside the
support $k^{(4)}(q)=C(144-1440q^2+1680q^4)$.
Its absolute integral is elementary, splitting at
$q^2=(15\pm2\sqrt{30})/35$.

\begin{corollary}[A growing stationary clock]\label{cor:clock}
Fix $\epsilon>0$, $Y>h>0$ and $\ell_0>0$. Choose
\begin{equation}
 N=\left\lfloor\frac{\mu\kappa^{1/3}}{\pi}\right\rfloor,\qquad
 X=N\pi/\mu\sim\kappa^{1/3}.
 \label{eq:growing}
\end{equation}
For all sufficiently large $t$, the finite gates hold and
\begin{equation}
 K*\rh(t,x,y)=\frac{L(t)}{\pi}k_h(y)
       +O_{Y,h,\ell_0,\epsilon}(\kappa^{-1/2}+L^{-3}),
 \label{eq:clock}
\end{equation}
uniformly under \eqref{eq:buffer} for $\ell\ge\ell_0$.
The width $\ell=\ell(t,x)$ may vary between observations.
\end{corollary}
\begin{proof}
The moment bounds imply $vZ^2=O(\kappa^{-1/3})$,
$\delta=O(\kappa^{-1/2})$ and $r=O(\kappa^{-1/2}/L)$.
The explicit upper allowances have these orders as well.
Thus $n_\ell r/\ell^2=O(\delta/\ell+\delta/(L\ell^2))$;
the phase term is $O(L^{-3}\ell^{-4})$.
Finally
$q_t(L)=tL^2-tL/2+9L-9/4$, so $\eta=\epsilon/M$ decays faster
than every inverse power of $\kappa$.
Using the explicit allowances in \eqref{eq:Ereg} makes that term
negligible at every fixed algebraic order, uniformly in the stated widths.
\end{proof}
In particular $\pi(K*\rh)/L\to k_h(y)$ with error
$O(\kappa^{-1/2}/L+L^{-4})$ on this growing domain.
Replacing $L$ by the elementary logarithmic expression in
\eqref{eq:Lambert} adds an absolute
$O(\log\log t/\log t)$ error.
A positive time average is obtained by averaging both
the model $L(t)k_h/\pi$ and the pointwise allowances over time.
It is not obtained by evaluating $L$ at the mean time.
